\vspace{-0.3em}\begin{center}{\large\bfseries Abstract}\end{center}
\vspace{-1.1em}
\begin{list}{}{\setlength{\leftmargin}{1.6em}\setlength{\rightmargin}{1.6em}\setlength{\topsep}{0pt}}\item\relax
\small
Short-video platforms are often said to trap users in ``information cocoons''. Existing measurements, and the cocoon self-checks that Chinese platforms launched in 2025, count how narrow a feed is across \emph{topics}. They cannot tell a user who simply likes a topic from a user who, within that topic, is shown one side of an argument more and more often. This gap matters for everyday issues such as marriage, appearance and career choices, where one topic can be told through very different narratives. We study this \emph{viewpoint cocoon} on the public Kuaishou RecIF logs (160{,}412 users, 75.3 million impressions) and the open OneRec recommender trained on them. We define viewpoint skew $V$ as the log-odds of meeting a narrative within an issue, relative to the platform average, and split total skew exactly into a platform layer and a personal layer. Language-model labels show that \nOppShare\% of issue impressions carry an opposing narrative and that such videos draw 1.9 times the comments. For gender group labels in marriage videos, platform ranking raises exposure by a factor of \nPlatformOdds{} over supply, and personalisation adds a further factor of \nOddsStrong{} for users who finish such videos, while their topic share stays the same. Detecting a doubling takes about \nDetNDouble{} issue videos per sample, close to four months of reference viewing. Of eleven escape strategies, topic-level controls remove the skew only by removing the topic, the ``not interested'' button does worse than scrolling past, and every gain fades within a week of stopping. We turn these results into a design for a viewpoint exposure report.
\end{list}
